\slajdid{09_1-s046}
\begin{frame}[label=09_1-s046]
Napon logičke nule će se dobiti kao rešenje jednačine
\[k_{n,L}(V_{DD}-V_O)(V_{DD}+2\Delta-V_O)=k_{n,D}V_O(2V_I-2V_{Tn,D}-V_O)\]
za $V_I=V_{OH}$
\[k_{n,L}(V_{DD}-V_O)(V_{DD}+2\Delta-V_O)=k_{n,D}V_O(2V_{OH}-2V_{Tn,D}-V_O)\]
Uz $V_{OL}$ malo
\[k_{n,L}(V_{DD}-V_O)(V_{DD}+2\Delta)=k_{n,D}V_O(2V_{OH}-2V_{Tn,D})\]
\[V_{OL}\approx\frac12\frac{k_{n,L}}{k_{n,D}}\frac{V_{DD}(V_{DD}+2\Delta)}{V_{DD}-V_{Tn,D}}=\frac12\frac{k_{n,L}}{k_{n,D}}\frac{V_{DD}(2V_C-V_{DD}-V_{Tn,L})}{V_{DD}-V_{Tn,D}}\]
\end{frame}
